\section{Conclusion and Future Work}

This project develops a lightweight sparse Python solver for onset-level Taylor-cone modeling. The solver couples an axisymmetric electrostatic solution to a Young--Laplace--Maxwell interface residual, reconstructs the interface field on a rounded cone, and recovers the classical Taylor half-angle. Its discretization is verified against manufactured solutions, against a staircase representation of the same curved boundary, and against an independent body-fitted finite-element discretization.

\rev{The verification study reported in Section~\ref{sec:completed-verification} covers an independent finite-element comparison, curved-boundary refinement, angle and parameter sensitivity, and a synthetic-image round trip of the extraction pipeline that meets its preset $1^\circ$ angle and 2\% profile acceptance limits over the two lowest noise levels. The public-image analysis of Section~\ref{sec:experimental-procedure} reports a mean absolute half-angle difference of $1.337^\circ$ between two published image-derived series and a systematic negative residual trend under a fixed-pressure closure. Because the extraction chain is qualified only for that numerical round trip and not against recorded images, the image-derived angles are presented as a published-image comparison rather than as a qualified measurement. The study therefore motivates the planned experimental validation rather than replacing it.}

Future work includes improving sharp-interface boundary representation, coupling threshold space charge inside the shape-optimization loop, adding more physically constrained charge/current closures, extending toward leaky-dielectric effects where justified, and comparing against higher-fidelity simulations or published experimental datasets. The intended contribution is not to replace full multiphysics simulation, but to provide a transparent, reproducible, and computationally accessible model for Taylor-cone onset studies.
